Un ascenseur avec ses occupants a une masse $\mathrm{m}=\qty{850}{\kg}$. Durant
un déplacement $AB$, \textbf{il monte à vitesse constante} sous l'effet de la
force de traction exercée par le câble tracteur.
\begin{enumerate}
    \item En négligeant les forces de frottement, faire le bilan des forces
          exercées sur l'ascenseur (nom, caractéristiques
          connues).~\textbf{(0,5pt)}
    \item Peut-on affirmer que le principe d'inertie s'applique au skieur~?
          Justifier votre réponse.~\textbf{(1pt)}
    \item En déduire la relation entre les forces exercées sur l'ascenseur et
          déterminer la valeur de la force de traction.~\textbf{(1pt)}
    \item L'ascenseur a parcouru la distance $AB=\qty{60.0}{\m}$ en une durée
          $\Delta t=\qty{0.5}{\min}$. Calculer la valeur du travail du poids de
          l'ascenseur et de ses occupants durant ce déplacement. Est-il moteur
          ou résistant? Justifier la réponse.~\textbf{(1pt)}
    \item Déterminer la puissance moyenne de la force de traction exercée par le
          câble durant le déplacement $AB$.~\textbf{(1pt)}
\end{enumerate}
\begin{donnees}
    Intensité de la pesanteur $\mathrm{g}=\qty{9.80}{\N\per\kg}$.
\end{donnees}
